\scanpage{28}
% L1–3
sur l'ens d'indices $\Delta_n$ pour les $\sigma_i\in G$] Si \deleted{\unsure{$z$}} $x_i\in D(\{i\}),y_j\in D(\{j\})$, on pose $x_i\unsure{\prec|<}y_j$ ssi $i<j$ et $\exists z\in D_*$ tel que \deleted{\unsure{$z$}} $x_i,y_j\leq z$.
% L4–5
\deleted{Cette relation d'ordre satisfait \textcircled{b} \inserted{(par définition)} et \textcircled{c} (par définition aussi)}
% L6–8
Cette relation est-elle transitive ? Supposons $i<j<k$, $x_i\in D(\{i\}),y_j\in D(\{j\}),z_k\in D(\{k\})$, \deleted{\unsure{$x_i$}} $x_i,y_j$ viennent de $u_{ij}\in D(\{i,j\})$, $y_j,z_k$ viennent
% L9–10
de $v_{jk}\in D(\{\deleted{\unsure{i}}j,k\}.$ Donc $(u_{ij},v_{jk})$ est un élément de
\[
D(\{i,j\})\mathbin{\times}_{D(\{j\})}D(\{j,k\})
\]
% L11–14
on serait content qu'il provienne d'un \unsure{élement|élmt|élément} $w_{ijk}\in D(\{i,j,k)\}$, car il \unsure{devrais|devrait|alors|dernier} définirait un élément \deleted{\unsure{$u$}} $\unsure{\chi|\gamma|x}_{ik}\in D(\{i,k\})$ \deleted{dont} se restreignant suivant $x_i$ et $\unsure{y_k|z_k}$. Donc
% L15
on veut
\[
D(\{i,j,k\})\xrightarrow{\text{\textcircled{épi}}}D(\{i,j\})\mathbin{\times}_{D(\{j\})}D(\{j,k\})
\]
% L16
Or je dis plus généralement qu'on a

% L17–21
\uline{Lemme} Soient $H,H'\subset\Delta_n$ tels que $H\cap H'\ne\varnothing$ et $H\cap H'$ soit un segment final de l'ens tot. ordonné $H$, un segment initial de l'ens tot ordonné $H'$. Considérons l'appl. canonique
% L22
\[
D(H\cup H')\xrightarrow{\deleted{\unsure{\sim}}}D(H)\mathbin{\times}_{\substack{\deleted{\unsure{D(H)\cap D(H')}}\\ D(H\cap H')}}D(H')
\]
% L23
Celle-ci est \uline{surjective}.
% archive sheet number: 27
